%==============================================================================
%  CAPÍTULO — EL MODELO ESTADOUNIDENSE
%==============================================================================
\chapter{El modelo estadounidense: el \emph{Bankruptcy Code}}
\label{cap:comparado-eeuu}

\fichaCapitulo{El ordenamiento que fijó el paradigma de la reorganización y del
\freshstart{}.}{Estructura del Título 11 del \emph{U.S.\ Code}: liquidación del Capítulo 7,
reorganización del Capítulo 11, plan de la persona natural del Capítulo 13 y reconocimiento
transfronterizo del Capítulo 15.}

%==============================================================================
\bloqueTeorico
%==============================================================================

\subsection{Por qué el derecho estadounidense es el comparador obligado}

Ningún ordenamiento ha influido tanto en el diseño concursal contemporáneo como el
estadounidense. Dos de sus construcciones se han universalizado: la \destacar{reorganización con
el deudor en posesión} ---que inspiró la reestructuración preventiva europea y, mediatamente, la
\LIR{} chilena--- y la \destacar{descarga de saldos insolutos} como derecho del deudor honesto, no
como gracia de los acreedores.

La razón de esa influencia es histórica y constitucional. La \emph{Bankruptcy Clause} del artículo
I, sección 8, de la Constitución de 1787 entregó al Congreso federal la potestad de dictar
\enquote{leyes uniformes en materia de bancarrota}, lo que produjo un estatuto único para toda la
Unión ---el \destacar{Título 11 del \emph{United States Code}}, o \emph{Bankruptcy Code}
\parencite{us-bankruptcy-code}--- aplicado por tribunales especializados: las \emph{bankruptcy
courts}, unidades de las cortes federales de distrito servidas por jueces con dedicación exclusiva
a la materia.

\begin{foco}[La lección institucional que Chile no ha tomado]
La especialización orgánica es el rasgo estadounidense más difícil de trasplantar y, quizá, el más
relevante. Chile entrega los concursos a juzgados civiles de competencia común, con la carga
general que ello implica; Estados Unidos los entrega a jueces que sólo ven insolvencia. Buena parte
de la calidad y la celeridad del sistema norteamericano se explica por ese dato orgánico, no por
sus normas sustantivas.
\end{foco}

\subsection{La arquitectura del Título 11}

El \emph{Code} se organiza en capítulos que no son procedimientos alternativos entre sí, sino vías
diferenciadas por el tipo de deudor y por la finalidad perseguida:

\begin{description}[leftmargin=0pt,style=nextline,font=\sffamily\bfseries]
  \item[Capítulos 1, 3 y 5 — Parte general] Definiciones, administración del caso y régimen del
  patrimonio concursal (\emph{estate}), acreedores y descarga. Se aplican transversalmente a los
  demás capítulos.
  \item[Capítulo 7 — Liquidación] Realización del activo por un \emph{trustee} y distribución
  entre acreedores. Es la vía por defecto.
  \item[Capítulo 9 — Municipalidades] Reorganización de entidades públicas locales; sin
  equivalente chileno.
  \item[Capítulo 11 — Reorganización] Reestructuración de la empresa, típicamente con el deudor
  conservando la administración.
  \item[Capítulo 12 — Agricultores y pescadores] Estatuto especial de reorganización.
  \item[Capítulo 13 — Plan de la persona con ingresos regulares] Reestructuración de deudas de la
  persona natural mediante un plan de pagos con cargo a rentas futuras.
  \item[Capítulo 15 — Insolvencia transfronteriza] Incorporación de la Ley Modelo
  \textsc{cnudmi}/\textsc{uncitral}, matriz del Capítulo VIII chileno.
\end{description}

\subsection{Los dos principios estructurantes}

\paragraph{El \emph{automatic stay} (§ 362).} La sola presentación de la solicitud produce, de
pleno derecho y sin resolución judicial que lo declare, la \destacar{paralización automática} de
toda acción de cobro, ejecución, embargo o perfeccionamiento de garantías contra el deudor y su
patrimonio. Es el equivalente funcional ---y el antecedente--- de la \pfc{} del \art{57} chileno,
con dos diferencias notables: opera \emph{ipso iure} desde la presentación, sin necesidad de
resolución, y no tiene plazo fijo de duración, sino que se levanta caso a caso mediante
\emph{motion for relief from stay}.

\paragraph{El \emph{discharge}.} La extinción de los saldos insolutos es el corazón del sistema y
se concibe como el derecho del \enquote{deudor honesto pero desafortunado} ---la fórmula clásica
de \emph{Local Loan Co.\ v.\ Hunt}--- a un nuevo comienzo. A diferencia de Chile, el \emph{Code}
contiene un \destacar{catálogo legal expreso de deudas no exonerables} en la § 523, cuya ausencia
en la \LIR{} explica buena parte de la controversia chilena sobre el CAE (capítulo
\ref{cap:discharge-cae}).

\begin{alerta}[La ausencia del catálogo del § 523 es el nudo chileno]
El § 523 excluye de la descarga, entre otras, las deudas alimenticias, las tributarias recientes,
las nacidas de dolo o fraude, las multas y sanciones, las indemnizaciones por daños causados en
estado de ebriedad y ---con la relevancia que se verá--- los \destacar{créditos educativos}. Que
Chile carezca de una lista equivalente obliga a los tribunales a construirla por vía
interpretativa, con el resultado de fallos contradictorios entre salas de una misma Corte.
\end{alerta}

%==============================================================================
\bloqueNormativo
%==============================================================================

\subsection{El Capítulo 7: liquidación y descarga}

El Capítulo 7 abre con la constitución del \emph{estate} ---universalidad de bienes del deudor a la
fecha de la solicitud--- y la designación de un \emph{trustee}, auxiliar equivalente al liquidador
chileno, encargado de reunir, realizar y distribuir el activo.

Dos instituciones lo singularizan frente al modelo chileno:

\paragraph{Las \emph{exemptions}.} El deudor persona natural conserva un conjunto de bienes exentos
de realización, cuyo catálogo puede elegir entre el federal (§ 522) y el del Estado de su
domicilio. Cubren típicamente un monto de la vivienda (\emph{homestead exemption}), el vehículo,
herramientas de trabajo, enseres y fondos previsionales. La lógica es la misma de los bienes
inembargables chilenos (\artcpc{445}), pero con una amplitud considerablemente mayor y una enorme
dispersión entre Estados.

\paragraph{El \emph{means test}.} Introducido por la reforma \textsc{bapcpa} de 2005, somete al
deudor persona natural a un examen de capacidad de pago: si su ingreso supera la mediana estatal y
puede destinar un excedente al pago de sus acreedores, su caso puede ser desestimado por
\emph{abuse} o reconducido al Capítulo 13. Es un filtro de \emph{merecimiento} económico que la
\LIR{} chilena no contempla ---nuestro control equivalente es de \emph{probidad}, no de capacidad:
el incidente de mala fe del \art{169 A}---.

\begin{foco}[Dos filtros distintos para el mismo problema]
Estados Unidos y Chile temen lo mismo ---el uso oportunista de la descarga--- pero lo filtran de
manera opuesta. El \emph{means test} pregunta \emph{¿puede usted pagar?} y es un test objetivo de
ingresos. El \art{169 A} pregunta \emph{¿se comportó usted honestamente?} y es un test de conducta.
El primero excluye al deudor solvente; el segundo, al deudor desleal. Son respuestas distintas a
preguntas distintas, y conviene no confundirlas al invocar el derecho comparado.
\end{foco}

\subsection{El Capítulo 11: la reorganización y el \emph{debtor in possession}}

El Capítulo 11 es la construcción de mayor irradiación mundial. Sus piezas centrales:

\paragraph{El \emph{debtor in possession} (DIP).} Salvo fraude o mala administración manifiesta
---casos en que se designa un \emph{trustee}---, el deudor \destacar{conserva la administración} de
su empresa y asume los deberes fiduciarios de un \emph{trustee} frente a la masa. Es el reverso del
desasimiento: se presume que quien mejor conoce el negocio es quien debe conducirlo durante la
reestructuración. La reorganización chilena adopta la misma premisa ---el deudor conserva la
administración bajo supervisión del veedor---, aunque con facultades de disposición mucho más
acotadas (\art{74}).

\paragraph{El \emph{exclusivity period}.} Durante los primeros 120 días (prorrogables), sólo el
deudor puede presentar un plan. Vencido el plazo, cualquier interesado puede proponer el suyo, lo
que introduce competencia entre planes ---mecanismo ausente en Chile, donde la propuesta es siempre
del deudor---.

\paragraph{El \emph{disclosure statement}.} Antes de votar, los acreedores deben recibir un
documento con \enquote{información adecuada} aprobado judicialmente. Es un deber de transparencia
informativa más intenso que el del \art{56} chileno, y su incumplimiento invalida la votación.

\paragraph{Clases, votación y \emph{cram-down}.} El plan agrupa los créditos en clases; cada clase
aprueba con dos tercios del monto y más de la mitad en número de los votantes. Si una clase
disidente rechaza el plan, el tribunal puede aun así confirmarlo ---\emph{cram-down}--- siempre que
el plan no discrimine injustamente y sea \emph{fair and equitable}, lo que en los créditos valistas
se traduce en la \destacar{\emph{absolute priority rule}}: ninguna clase inferior puede recibir
nada mientras una clase superior disidente no haya sido pagada íntegramente. A ello se suma el
\emph{best interest of creditors test}: ningún acreedor puede recibir menos de lo que obtendría en
una liquidación del Capítulo 7.

\begin{alerta}[El arrastre de clases que Chile no tiene]
La \LIR{} carece de un mecanismo de \emph{cram-down} entre clases: si una clase rechaza la
propuesta, el acuerdo simplemente fracasa y sobreviene la liquidación (\art{96}). El derecho
estadounidense ---y, tras la Directiva \parencite{ue-directiva-2019-1023}, el europeo--- permite
imponer el plan a la clase disidente bajo condiciones de equidad. Es, probablemente, la reforma
pendiente de mayor calado para la reorganización chilena, y el debate sobre las quitas al
garantizado disidente (capítulo \ref{cap:reorganizacion-ordinaria}) es su antesala.
\end{alerta}

\paragraph{Las ventas del § 363.} El deudor puede enajenar activos fuera del giro ordinario, con
autorización judicial y \destacar{libres de gravámenes} (\emph{free and clear}), trasladándose las
garantías al producto de la venta. Esta técnica ---la \emph{363 sale}--- permitió reestructuraciones
célebres mediante la venta rápida de la empresa en marcha, y es el antecedente funcional de la
venta como unidad económica del derecho chileno.

\paragraph{El \emph{DIP financing} (§ 364).} Quien financia a la empresa durante el procedimiento
obtiene prioridad administrativa y, con autorización judicial, puede incluso obtener una garantía
de rango superior a las preexistentes (\emph{priming lien}). El \art{74} chileno recoge la idea
matriz ---preferencia del \artcc{2472} \Nro 4 para el nuevo financiamiento--- pero sin la
posibilidad de posponer garantías ya constituidas.

\paragraph{El \emph{Subchapter V}.} La \emph{Small Business Reorganization Act} de 2019 creó un
procedimiento simplificado para pequeñas empresas: sin comité de acreedores, sin
\emph{disclosure statement} formal, con plazos breves y sin aplicación estricta de la
\emph{absolute priority rule}. Su parentesco funcional con la reorganización simplificada chilena
de la \Lreforma{} es evidente, y ambas responden al mismo diagnóstico: el procedimiento ordinario
resultaba inalcanzable para la empresa pequeña.

\subsection{El Capítulo 13: el plan de la persona natural}

Reservado a personas naturales con ingresos regulares y deudas bajo ciertos umbrales, el Capítulo
13 permite conservar los bienes ---incluida la vivienda--- a cambio de destinar el ingreso
disponible a un plan de pagos de tres a cinco años, supervisado por un \emph{trustee}. Cumplido el
plan, opera la descarga del saldo.

Su interés comparado para Chile es directo: es el modelo funcional de la
\destacar{renegociación de la persona deudora} de los \art{260} y siguientes, aunque con dos
diferencias estructurales. La chilena es \destacar{administrativa y gratuita} ---se tramita ante la
\superir{}, no ante un tribunal--- y su horizonte temporal es mucho más breve. La estadounidense es
judicial y se extiende por años, pero permite algo que la chilena no: la \emph{cura} de la mora
hipotecaria para evitar la ejecución de la vivienda.

\subsection{Las deudas educativas y el \emph{undue hardship}}

El § 523(a)(8) excluye de la descarga los créditos educativos, salvo que el deudor acredite que su
pago le impondría un \destacar{\emph{undue hardship}}. Los tribunales han precisado el estándar
mediante el llamado \emph{test de Brunner}, de tres elementos copulativos: (i) que el deudor no
pueda mantener un nivel de vida mínimo si paga; (ii) que esa situación probablemente persista
durante una porción sustancial del período de pago; y (iii) que el deudor haya hecho esfuerzos de
buena fe por pagar.

\begin{practica}[El argumento comparado en la defensa del CAE]
La comparación es doblemente útil ante un tribunal chileno. Primero, demuestra que \emph{cuando un
legislador quiere excluir el crédito educativo de la descarga, lo dice expresamente}: el silencio
de la \LIR{} no puede suplirse por analogía con la Ley \Nro 20.027. Segundo, muestra que incluso el
ordenamiento que sí lo excluye admite una válvula de escape por \emph{hardship}, de modo que la
exclusión chilena ---absoluta e incondicionada--- resulta más severa que su modelo. Ambos
argumentos son utilizables en la línea de defensa del capítulo \ref{cap:discharge-cae}.
\end{practica}

\subsection{Las acciones de reintegración: \emph{preferences} y \emph{fraudulent transfers}}

El § 547 permite revocar las \destacar{\emph{preferential transfers}}: pagos o garantías otorgados
a un acreedor dentro de los 90 días anteriores a la solicitud ---o un año si se trata de un
\emph{insider}--- que le permitan recibir más de lo que habría obtenido en el Capítulo 7. Es una
acción \destacar{puramente objetiva}: no exige mala fe ni conocimiento, sólo el efecto preferencial
dentro del período. El paralelismo con la revocatoria objetiva del \art{287} chileno es estrecho,
como lo es el del plazo diferenciado para las personas relacionadas.

El § 548, por su parte, sanciona las \emph{fraudulent transfers}: enajenaciones con ánimo de
defraudar o, en su variante constructiva, realizadas sin recibir un valor razonablemente
equivalente estando el deudor insolvente. Equivale a la revocatoria subjetiva del \art{288},
aunque la modalidad \emph{constructive} ---que prescinde por completo del elemento intencional y se
satisface con el desequilibrio objetivo de la contraprestación--- carece de correlato exacto en el
derecho chileno.

\subsection{El Capítulo 15 y el reconocimiento transfronterizo}

El Capítulo 15 incorporó al derecho estadounidense la Ley Modelo \textsc{cnudmi}, la misma matriz
del Capítulo VIII chileno. Sus categorías son, por tanto, las que ya se examinaron: procedimiento
extranjero \emph{principal} o \emph{no principal} según el \textsc{comi}, acceso directo del
representante extranjero, efectos automáticos del reconocimiento del principal y deber de
cooperación entre tribunales.

Su relevancia práctica para Chile es concreta: los grandes casos chilenos de insolvencia
transfronteriza ---señaladamente el de \textsc{latam}--- se articularon como reconocimiento en
Chile de un procedimiento del Capítulo 11 tramitado en Nueva York (capítulo
\ref{cap:transfronteriza}). Conocer la mecánica estadounidense no es un lujo académico para el
abogado concursalista chileno: es el presupuesto para litigar la posición de los acreedores locales
en esos procedimientos.

%==============================================================================
\bloquePractico
%==============================================================================

\subsection{Tabla de equivalencias funcionales}

\begin{table}[htbp]
\centering
\small
\caption{Correspondencias entre el \emph{Bankruptcy Code} y la Ley \Nro 20.720}
\label{tab:cap25-equivalencias}
\begin{tabularx}{\textwidth}{@{}>{\raggedright\arraybackslash}p{0.30\textwidth}
                                 >{\raggedright\arraybackslash}p{0.26\textwidth} X@{}}
\toprule
\textbf{Estados Unidos} & \textbf{Chile} & \textbf{Diferencia relevante} \\
\midrule
\emph{Automatic stay} (§ 362) & \pfc{} (\art{57}) & En EE.UU.\ opera de pleno derecho desde la
presentación y sin plazo fijo; en Chile requiere resolución y tiene duración tasada \\
\addlinespace
\emph{Debtor in possession} (Cap.\ 11) & Deudor bajo supervisión del veedor & Facultades de
disposición mucho más amplias en EE.UU. \\
\addlinespace
\emph{Cram-down} y \emph{absolute priority rule} & \emph{Sin equivalente} & Chile carece de arrastre
de clases disidentes: el rechazo conduce a liquidación \\
\addlinespace
\emph{363 sale} (libre de gravámenes) & Venta como unidad económica & En EE.UU.\ la garantía se
traslada al precio por regla general \\
\addlinespace
\emph{DIP financing} (§ 364) & Financiamiento del \art{74} & EE.UU.\ admite \emph{priming lien}
sobre garantías previas; Chile no \\
\addlinespace
\emph{Subchapter V} & Reorganización simplificada (\art{286 A} y ss.) & Misma finalidad: acceso de
la empresa pequeña \\
\addlinespace
Capítulo 13 & Renegociación de la persona deudora & Chile: administrativa y gratuita; EE.UU.:
judicial y de 3 a 5 años \\
\addlinespace
\emph{Exemptions} (§ 522) & Bienes inembargables & Catálogo mucho más amplio y variable por Estado \\
\addlinespace
\emph{Means test} & Incidente de mala fe (\art{169 A}) & Filtro de capacidad de pago \emph{vs.}
filtro de probidad \\
\addlinespace
Excepciones del § 523 & \emph{Sin catálogo equivalente} & Origen de la controversia chilena sobre
el CAE \\
\addlinespace
\emph{Preferences} (§ 547) & Revocatoria objetiva (\art{287}) & Ambas objetivas, con plazo agravado
para \emph{insiders}/personas relacionadas \\
\addlinespace
Capítulo 15 & Capítulo VIII & Misma matriz: Ley Modelo \textsc{cnudmi} \\
\bottomrule
\end{tabularx}
\end{table}

\begin{foco}[Qué tomar y qué no del modelo estadounidense]
Tres instituciones estadounidenses son trasplantables y están pendientes en Chile: el
\destacar{catálogo legal de deudas no exonerables}, que cerraría la controversia del CAE; el
\destacar{arrastre de clases disidentes} con sus contrapesos de equidad; y la
\destacar{especialización orgánica} de los tribunales. En cambio, el \emph{means test} responde a
una filosofía distinta de la chilena ---mérito económico frente a probidad--- y su importación
descontextualizada endurecería el acceso al procedimiento sin atacar el problema que la \Lreforma{}
quiso resolver.
\end{foco}
